\documentclass[10pt,a4paper]{amsart}
\usepackage[margin=19mm]{geometry}
\usepackage{amsmath,amssymb,mathtools}
\usepackage[scr=rsfs,cal=euler]{mathalfa}
\usepackage{xfrac}
\usepackage[unicode,hidelinks,bookmarks=false]{hyperref}
\newcommand{\ie}{i.e.\ }
\newcommand{\cf}{cf.\ }
\newcommand{\resp}{respectively\ }
\newcommand{\Subsection}{\subsection}
\providecommand{\nobreakdash}{\nobreak\hbox{-}}
\setlength{\emergencystretch}{2em}
\multlinegap0pt
\usepackage{algorithm}\usepackage{algpseudocode}
\usepackage{amssymb}
\usepackage{braket}
\usepackage[shortlabels]{enumitem}
\usepackage{mathtools}
\usepackage{tikz}
\usepackage{xcolor}
\newtheorem{lemma}{Lemma}
\newcommand{\R} {\mathbb R}
\title{Identifying Drift, Diffusion, and Causal Structure from Temporal Snapshots}
\author{Vincent Guan, Joseph Janssen, Hossein Rahmani, Andrew Warren, Stephen Zhang, Elina Robeva, Geoffrey Schiebinger}
\date{Source: arXiv:2410.22729v4 (CC BY 4.0)}
\begin{document}
\maketitle

\noindent\emph{Source and licence.} Identifying Drift, Diffusion, and Causal Structure from Temporal Snapshots. Version arXiv:2410.22729v4 (CC BY 4.0), distributed by arXiv under the Creative Commons Attribution 4.0 International licence, http://creativecommons.org/licenses/by/4.0/, as recorded in the arXiv licence metadata for this exact version and shown on the abstract page https://arxiv.org/abs/2410.22729v4. Full licence notice: \texttt{licenses/arxiv-2410.22729-CC-BY-4.0.txt}.\par\smallskip

\noindent\textit{Editorial scope.} Counted unit: the article's lemma lem:conjugate\_rotation (source0.tex lines 1548--1559). The article's complete original proof of it is delivered with it (source0.tex lines 1561--1581, one \texttt{proof} environment of 21 lines). No mathematics is written, repaired or re-derived: the statement and the proof are the article's own lines, in the article's own notation, and no retained line is substituted or re-flowed.  No further source-local context is needed: every cross-reference of every retained line resolves inside the two delivered spans.  Nothing was omitted from any retained span, so the package carries no omission record.  No retained line contains a citation, so the wrapper carries no bibliography and no bibliography entry.  No retained line contains a figure environment, an image-inclusion command or drawing code, so no image is delivered and the package declares no asset dependency.  Editorial formatting only, in addition: an AMS \texttt{amsart} preamble that restates the article's own declarations and macros used by the retained lines, namely the environments \texttt{lemma} on separate counters, together with the standard packages those lines need.  The retained environments print in this document as Lemma 1 in order of appearance, whereas the article prints the same items with the numbering of its own full document.  The complete native source of the article as distributed in the arXiv e-print for this exact version is bundled unchanged as \texttt{sources/167660/source0.tex}.
\begin{lemma}
Let \(\Sigma \succ 0\) and \(A \in \R^{d \times d}\) be such that
\[
A\Sigma + \Sigma A^\top  = 0.
\]
Then, $B = \Sigma^{1/2} A\,\Sigma^{-1/2}$  is skew-symmetric, and for all \(t\in\R\),
\[
e^{At} = \Sigma^{-1/2}\,e^{B\theta}\,\Sigma^{1/2}.
\]
In other words, the $\Sigma$-generalized rotation \(e^{A\theta}\) produces a classical rotation \(e^{B\theta}\) after the change of variables \(y = \Sigma^{1/2}x\), which can be interpreted as a shearing and scaling transformation of $\R^d$.
\label{lem:conjugate_rotation}
\end{lemma}

\begin{proof}
Since \(\Sigma \succ 0\), the unique symmetric positive definite square root \(\Sigma^{1/2}\) exists and is invertible. Define
\[
B = \Sigma^{-1/2} A\,\Sigma^{1/2}, \qquad B^\top =  \Sigma^{1/2} A^\top\,\Sigma^{-1/2}.
\]
Then, since \(A\Sigma + \Sigma A^\top  = 0\), we can multiply on the left by \(\Sigma^{-1/2}\) and on the right by \(\Sigma^{-1/2}\) to yield
\[
0 = \Sigma^{-1/2} A\,\Sigma^{1/2} + \Sigma^{1/2}\,A^\top \,\Sigma^{-1/2} = B + B^\top,
\]
which proves that \(B\) is skew-symmetric, and therefore \(e^{B\theta}\) defines a classical rotation in \(\R^d\), such that
\[
e^{B\theta} = e^{\Sigma^{-1/2}A\Sigma^{1/2}}=\Sigma^{-1/2}e^{A\theta}\Sigma^{1/2},
\]
where the last equality can be seen by writing $e^{B\theta} = \sum_{n=0}^{\infty} \theta^n \frac{B^n}{n!}$ and noting that $B^n = (\Sigma^{-1/2}A\Sigma^{1/2})^n = \Sigma^{-1/2}A^n\Sigma^{1/2}$. Rearranging yields $e^{A\theta} = \Sigma^{-1/2}\,e^{B\theta}\,\Sigma^{1/2}$ as desired.

To explicitly see that the change of variable \(y = \Sigma^{1/2} x\) amounts to a shearing and scaling, note that any positive definite matrix admits an \(LDL^\top \) factorization:
\[
\Sigma = L\,D\,L^\top ,
\]
where \(L\) is unit lower-triangular (so its off-diagonal entries describe a shear transformation) and \(D\) is a diagonal matrix with positive entries (representing scaling along the coordinate axes). Thus, the transformation \(x\mapsto y=\Sigma^{1/2} x\) decomposes into a shear (via \(L\)) and a scaling (via \(D^{1/2}\)). 
\end{proof}

\end{document}
